\documentclass[a4paper]{article}

% \usepackage[default]{fontsetup}

\usepackage{fancyhdr}
\usepackage{extramarks}
\usepackage{amsmath}
\usepackage{amsthm}
\usepackage{amsfonts}
\usepackage{tikz}
\usepackage[plain]{algorithm}
\usepackage{algpseudocode}
\usepackage{enumerate}
\usepackage{tikz}

\usetikzlibrary{automata,positioning}

%
% Basic Document Settings
%  

\topmargin=-0.2in
\evensidemargin=0in
\oddsidemargin=0in
\textwidth=6.5in
\textheight=9.5in
\headsep=0.25in

\linespread{1.1}

\pagestyle{fancy}
\lhead{\hmwkAuthorName}
\chead{\hmwkClass : \hmwkTitle}
\rhead{\firstxmark}
\lfoot{\lastxmark}
\cfoot{\thepage}

\renewcommand\headrulewidth{0.4pt}
\renewcommand\footrulewidth{0.4pt}

\setlength\parindent{0pt}

%
% Create Problem Sections
%

\newcommand{\enterProblemHeader}[1]{
    \nobreak\extramarks{}{Problem \arabic{#1} continued on next page\ldots}\nobreak{}
    \nobreak\extramarks{Problem \arabic{#1} (continued)}{Problem \arabic{#1} continued on next page\ldots}\nobreak{}
}

\newcommand{\exitProblemHeader}[1]{
    \nobreak\extramarks{Problem \arabic{#1} (continued)}{Problem \arabic{#1} continued on next page\ldots}\nobreak{}
    \stepcounter{#1}
    \nobreak\extramarks{Problem \arabic{#1}}{}\nobreak{}
}

\newcommand*\circled[1]{\tikz[baseline=(char.base)]{
		\node[shape=circle,draw,inner sep=2pt] (char) {#1};}}


\setcounter{secnumdepth}{0}
\newcounter{partCounter}
\newcounter{homeworkProblemCounter}
\setcounter{homeworkProblemCounter}{1}
\nobreak\extramarks{Problem \arabic{homeworkProblemCounter}}{}\nobreak{}

%
% Homework Problem Environment
%
% This environment takes an optional argument. When given, it will adjust the
% problem counter. This is useful for when the problems given for your
% assignment aren't sequential. See the last 3 problems of this template for an
% example.
%

\newenvironment{homeworkProblem}[2][-1]{%
    \ifnum#1>0\relax
        \setcounter{homeworkProblemCounter}{#1}
    \fi
    \section{Problem \arabic{homeworkProblemCounter} (#2 points)}
    \setcounter{partCounter}{1}
    \enterProblemHeader{homeworkProblemCounter}
}{%
    \exitProblemHeader{homeworkProblemCounter}
}

%
% Homework Details
%   - Title
%   - Class
%   - Due date
%   - Name
%   - Student ID

\newcommand{\hmwkTitle}{Homework\ \#07}
\newcommand{\hmwkClass}{Probability \& Statistics for EECS}
\newcommand{\hmwkDueDate}{2026-11-8}
\newcommand{\hmwkAuthorName}{}
\newcommand{\hmwkAuthorID}{}


%
% Title Page
%

\title{
    \vspace{2in}
    \textmd{\textbf{\hmwkClass:\\  \hmwkTitle}}\\
    \normalsize\vspace{0.1in}\small{Due\ on\ \hmwkDueDate\ at 23:59}\\
	\vspace{4in}
}

\author{
	Name: Zhang Yiheng\textbf{\hmwkAuthorName} \\
	Student ID: 2025531068\hmwkAuthorID}
\date{}

\renewcommand{\part}[1]{\textbf{\large Part \Alph{partCounter}}\stepcounter{partCounter}\\}

%
% Various Helper Commands
%

% Useful for algorithms
\newcommand{\alg}[1]{\textsc{\bfseries \footnotesize #1}}
% For derivatives
\newcommand{\deriv}[1]{\frac{\mathrm{d}}{\mathrm{d}x} (#1)}
% For partial derivatives
\newcommand{\pderiv}[2]{\frac{\partial}{\partial #1} (#2)}
% Integral dx
\newcommand{\dx}{\mathrm{d}x}
% Alias for the Solution section header
\newcommand{\solution}{\textbf{\large Solution}}
% Probability commands: Expectation, Variance, Covariance, Bias
\newcommand{\E}{\mathrm{E}}
\newcommand{\Var}{\mathrm{Var}}
\newcommand{\Cov}{\mathrm{Cov}}
\newcommand{\Bias}{\mathrm{Bias}}

\begin{document}


% \maketitle
% \thispagestyle{empty}
% \pagebreak

\date{
Due on Nov. 8, 2026, 23:59 UTC+8}
\title{SI 140A.02  Probability \& Statistics for EECS, Fall 2026 \\
Homework 7}
\maketitle
Read all the instructions below carefully before you start working on the assignment, and before you make a submission.
\begin{itemize}
    \item You are required to submit your homework in PDF format generated by \textbf{LaTeX}, otherwise you will get no points of this homework.
    \item You are required to write down all the major steps towards making your conclusions; otherwise you may obtain limited points of the problem.
    \item Write your homework in English; otherwise you will get no points of this homework.
    \item Any form of plagiarism will lead to $0$ point of this homework. 
\end{itemize}
\newpage
% Adapted from Joseph, Ch. 4, #24
\begin{homeworkProblem}{20}
Calvin and Hobbes play independent games. Calvin wins each game with probability $p$, and Hobbes wins with probability $q=1-p$, where $0<p<1$. The match ends as soon as either player has won two more games than the other.
\begin{enumerate}
\item[(a)] Let $T$ be the total number of games played. Find $\Pr(T=2k)$, $k=1,2,\ldots$. \emph{Hint:} Consider the games in consecutive pairs.
\item[(b)] Find $\E[T]$. Which value of $p$ maximizes this expectation?
\item[(c)] Find the probability that Calvin wins the match.
\end{enumerate}

\solution
\begin{enumerate}[(a)]
\item At the start of a pair, the scores are tied. A split pair (CH or HC)
has probability $2pq$ and restores the tie without ending the match;
a pair CC or HH ends the match, with total probability $p^2+q^2$.
Thus the match can only end after an even number of games, and independence gives
\[
\boxed{\Pr(T=2k)=(2pq)^{k-1}(p^2+q^2),\qquad k=1,2,\ldots.}
\]
\item The number of pairs $T/2$ is geometric on $\{1,2,\ldots\}$
with success probability $p^2+q^2$. Hence
\[
\boxed{\E[T]=\frac{2}{p^2+q^2}
=\frac{2}{2(p-\tfrac12)^2+\tfrac12}.}
\]
The denominator is minimized uniquely at $\boxed{p=1/2}$,
where the maximum expected duration is $\boxed{4}$ games.
\item Calvin wins exactly when some number of split pairs is followed by CC.
Summing these disjoint possibilities gives
\[
\boxed{\Pr(\text{Calvin wins})
=\sum_{k=1}^{\infty}(2pq)^{k-1}p^2
=\frac{p^2}{1-2pq}=\frac{p^2}{p^2+q^2}.}
\]
Here $2pq\le 1/2<1$, so the geometric series converges and the match
ends with probability one.
\end{enumerate}

\end{homeworkProblem}
\newpage
% Adapted from Joseph, Ch. 4, #40
\begin{homeworkProblem}{20}
A box contains $100$ shoelaces, each with two loose ends. Repeatedly pick two uniformly random loose ends and tie them together. A tie either joins two open pieces or closes one piece into a loop. Continue until all ends have been tied. Let $L$ be the total number of loops formed.
\begin{enumerate}
\item[(a)] When $m$ open pieces remain, find the probability that the next tie creates a loop.
\item[(b)] Find $\E[L]$ as a finite sum. Independence of the loop indicators is \emph{not} assumed.
\item[(c)] Find the probability that only one loop is formed by the end. A finite product is acceptable.
\end{enumerate}

\solution
\begin{enumerate}[(a)]
\item With $m$ open pieces, there are $2m$ loose ends. Of the
$\binom{2m}{2}$ equally likely pairs, exactly $m$ consist of the two
ends of the same piece. Therefore, conditional on any previous history,
\[
\boxed{\Pr(\text{next tie creates a loop}\mid m\text{ open pieces})
=\frac{m}{\binom{2m}{2}}=\frac{1}{2m-1}.}
\]
\item Each tie reduces the number of open pieces by one, whether it
joins two pieces or closes one into a loop. Thus the process visits
$m=100,99,\ldots,1$ exactly once. Let $J_m$ indicate that a loop is
created at the step with $m$ open pieces. By linearity of expectation,
\[
\boxed{\E[L]=\E\!\left[\sum_{m=1}^{100}J_m\right]
=\sum_{m=1}^{100}\frac{1}{2m-1}.}
\]
No independence of the indicators is needed.
\item At $m=1$, a loop is certain to form. Consequently, $L=1$ if and
only if no loop is created at any of the preceding steps. Applying the
chain rule and the history-independent conditional probability in (a),
\[
\boxed{\Pr(L=1)=\prod_{m=2}^{100}\left(1-\frac{1}{2m-1}\right)
=\prod_{m=2}^{100}\frac{2m-2}{2m-1}
=\frac{198!!}{199!!}.}
\]
\end{enumerate}

\end{homeworkProblem}
\newpage
% Adapted from Joseph, Ch. 4, #28 (zero-inflated Poisson)
\begin{homeworkProblem}{20}
A monitoring device is \emph{inactive} on a given day with probability $\alpha\in(0,1)$. On an active day, each of $n$ independent sensors produces an alert with a small probability $u$. Approximate the active-day alert count by $\mathrm{Pois}(\lambda)$, where $\lambda=nu$ (with large $n$ and small $u$). Let $X$ be the number of alerts on a randomly selected day. The inactive/active status is independent of the sensor outcomes.
\begin{enumerate}
\item[(a)] Identify the original distribution of the active-day alert count before approximation. Under the Poisson approximation, find $\Pr(X=0)$ and $\Pr(X=k)$ for all integers $k\ge1$.
\item[(b)] Find $\E[X]$ and $\Var(X)$. You may introduce independent variables $I\sim\mathrm{Bern}(\alpha)$ and $Y\sim\mathrm{Pois}(\lambda)$ with $X=(1-I)Y$.
\item[(c)] Take $n=200$, $u=0.01$ and $\alpha=1/5$ (so $\lambda=2$). Compare $\Pr(X=0)$ under the mixture model with the zero probability under $\mathrm{Pois}(\E[X])$. Explain briefly why using only the mean to select a Poisson model can be misleading.
\end{enumerate}

\solution
\begin{enumerate}[(a)]
\item Before approximation, the active-day alert count has distribution
$\mathrm{Bin}(n,u)$. Under the Poisson approximation, conditioning on the
device status gives the zero-inflated Poisson PMF
\[
\boxed{\Pr(X=0)=\alpha+(1-\alpha)e^{-\lambda},\qquad
\Pr(X=k)=(1-\alpha)e^{-\lambda}\frac{\lambda^k}{k!},\quad k\ge1.}
\]
\item Let $I\sim\mathrm{Bern}(\alpha)$ and
$Y\sim\mathrm{Pois}(\lambda)$ be independent, with $X=(1-I)Y$.
Since $\E[Y]=\lambda$ and $\E[Y^2]=\lambda+\lambda^2$,
\begin{align*}
\E[X]&=\E[1-I]\E[Y]=(1-\alpha)\lambda,\\
\E[X^2]&=\E[(1-I)^2]\E[Y^2]
=(1-\alpha)(\lambda+\lambda^2).
\end{align*}
Subtracting the square of the mean yields
\[
\boxed{\E[X]=(1-\alpha)\lambda,\qquad
\Var(X)=(1-\alpha)\lambda+\alpha(1-\alpha)\lambda^2.}
\]
These moments refer to the Poisson approximation in (a).
\item Here $\E[X]=(4/5)\cdot2=8/5$. The two zero probabilities are
\[
\boxed{\Pr(X=0)=\frac15+\frac45e^{-2}\approx0.30827,\qquad
\Pr\!\left(\mathrm{Pois}(8/5)=0\right)=e^{-8/5}\approx0.20190.}
\]
Inactivity creates additional zeros. More generally, strict convexity of
$e^{-x}$ gives
$\alpha+(1-\alpha)e^{-\lambda}>e^{-(1-\alpha)\lambda}$ for
$\lambda>0$. The mixture also has variance $2.24$, larger than its mean
$1.6$, whereas a Poisson variable has equal mean and variance.
Thus matching only the mean misses both the excess zeros and the extra variability.
\end{enumerate}

\end{homeworkProblem}
\newpage
% Original: overlapping patterns and second moments
\begin{homeworkProblem}{20}
A fair coin is tossed independently $n\ge5$ times. Let $X$ count the occurrences of the pattern HTH in consecutive positions, allowing overlaps (e.g., HTHTH contains two occurrences). For $i=1,\ldots,n-2$, let $I_i$ indicate the occurrence starting at position $i$, so $X=\sum_{i=1}^{n-2}I_i$.
\begin{enumerate}
\item[(a)] Find $\E[X]$.
\item[(b)] Compute $\Pr(I_i=I_{i+1}=1)$ and $\Pr(I_i=I_{i+2}=1)$. Explain why $I_i$ and $I_j$ are independent whenever $|i-j|\ge3$.
\item[(c)] Derive $\Var(X)$ as a function of $n$. Evaluate both $\E[X]$ and $\Var(X)$ when $n=20$.
\end{enumerate}

\solution
\begin{enumerate}[(a)]
\item Each length-three window equals HTH with probability $1/8$.
Linearity of expectation therefore gives
\[
\boxed{\E[X]=\sum_{i=1}^{n-2}\E[I_i]=\frac{n-2}{8}.}
\]
\item Occurrences starting one position apart are incompatible: the
shared toss at position $i+1$ would have to be both T and H. Occurrences
starting two positions apart require the five-toss string HTHTH. Thus
\[
\boxed{\Pr(I_i=I_{i+1}=1)=0,\qquad
\Pr(I_i=I_{i+2}=1)=\frac{1}{2^5}=\frac{1}{32}.}
\]
For $|i-j|\ge3$, the two indicators depend on disjoint sets of
independent tosses, so they are independent.
\item Each indicator has variance $(1/8)(7/8)=7/64$. Using (b),
\[
\Cov(I_i,I_{i+1})=-\frac{1}{64},\qquad
\Cov(I_i,I_{i+2})=\frac{1}{32}-\frac{1}{64}=\frac{1}{64},
\]
and all covariances at larger separations vanish. Among the $n-2$
starting positions, there are $n-3$ unordered pairs at distance one
and $n-4$ at distance two. Hence
\begin{align*}
\Var(X)&=\sum_{i=1}^{n-2}\Var(I_i)
+2\sum_{1\le i<j\le n-2}\Cov(I_i,I_j)\\
&=\frac{7(n-2)-2(n-3)+2(n-4)}{64}
=\boxed{\frac{7n-16}{64}}.
\end{align*}
For $n=20$, this gives
\[
\boxed{\E[X]=\frac94=2.25,\qquad \Var(X)=\frac{31}{16}=1.9375.}
\]
\end{enumerate}

\end{homeworkProblem}
\newpage
% Original: negative binomial, binomial, PGF and moments
\begin{homeworkProblem}{20}
A transmitter repeats independent attempts, each successful with probability $p\in(0,1)$. It stops at the \emph{third} successful attempt. Let $N$ be its total number of attempts, and $Y=N-3$ the number of failures before the third success.
\begin{enumerate}
\item[(a)] Find the PMF of $Y$ and identify its distribution.
\item[(b)] Express $\Pr(N\le6)$ in terms of a binomial random variable counting successes in the first six attempts. Evaluate this probability when $p=1/2$.
\item[(c)] View $Y$ as a sum of three independent geometric random variables, each counting failures before a success. Find its PGF $G_Y(t)=\E[t^Y]$ and use $G_Y'(1)$ and $G_Y''(1)$ to obtain $\E[N]$ and $\Var(N)$.
\end{enumerate}

\solution
\begin{enumerate}[(a)]
\item Write $q=1-p$. The event $Y=y$ requires exactly two successes
and $y$ failures in the first $y+2$ attempts, followed by a success.
Therefore
\[
\boxed{\Pr(Y=y)=\binom{y+2}{2}p^3q^y,\qquad y=0,1,2,\ldots.}
\]
Thus $Y$ is negative binomial with parameters $r=3$ and $p$, under
the convention that it counts failures before the third success.
\item Extend the independent attempt sequence beyond the stopping time,
and let $B$ count successes in its first six attempts. Then
$B\sim\mathrm{Bin}(6,p)$, and $N\le6$ if and only if $B\ge3$. Hence
\[
\boxed{\Pr(N\le6)=\Pr(B\ge3)
=\sum_{j=3}^{6}\binom{6}{j}p^j q^{6-j}.}
\]
When $p=1/2$, the probability is
\[
\boxed{\frac{20+15+6+1}{64}=\frac{21}{32}.}
\]
\item Let $Z_1,Z_2,Z_3$ count failures before the first success,
between the first and second successes, and between the second and third
successes, respectively. Independence of the attempts gives independent
$Z_j$ with $\Pr(Z_j=k)=pq^k$ for $k\ge0$, and $Y=Z_1+Z_2+Z_3$.
Consequently,
\[
G_{Z_j}(t)=\sum_{k=0}^{\infty}pq^k t^k=\frac{p}{1-qt},\qquad
\boxed{G_Y(t)=\left(\frac{p}{1-qt}\right)^3},\quad |qt|<1.
\]
Differentiating and evaluating at $t=1$ gives
\[
G_Y'(1)=\frac{3q}{p},\qquad G_Y''(1)=\frac{12q^2}{p^2}.
\]
Since $G_Y'(1)=\E[Y]$ and $G_Y''(1)=\E[Y(Y-1)]$,
\begin{align*}
\E[N]&=3+G_Y'(1)=\boxed{\frac{3}{p}},\\
\Var(N)&=\Var(Y)=G_Y''(1)+G_Y'(1)-[G_Y'(1)]^2\\
&=\frac{12q^2}{p^2}+\frac{3q}{p}-\frac{9q^2}{p^2}
=\boxed{\frac{3(1-p)}{p^2}}.
\end{align*}
\end{enumerate}

\end{homeworkProblem}
\end{document}
